
        In this section, we move towards looking at the empirical effects of the arrival of immigrants in the United States. In order to study the causal effects of Immigration in the United States, I am interested in estimating the following equation:
   \begin{equation}\label{eq:main}
     Y_{d,t} \;=\; \delta_t + \delta_{s(d)} + \beta \cdot \text{Immigration}_{d,t} + \varepsilon_{d,t}
    \end{equation}
    where $Y_{d,t}$ is a county level measure of business of county $d$ from $t-1$ to $t$. $I_{d,t}$ is the immigration to the county from time $t-1$ to $t$, $\delta_{s(i)}$ are state fixed-effects for county $d$, and $\delta_t$ are time fixed effects. The coefficient of interest is $\beta$. However, this simple equation is likely to suffer from endogeneity for a variety of reasons, such as reverse causality. The decision of where in the US to migrate is an endogenous choice, immigrants may choose to live in places with higher business dynamism. In order to extract causal estimates, I follow the IV approach outlined in \cite{ImmigrationInnovationGrowth}. I provide an elaborate explanation of the identification strategy in Section \ref{identification_appendix}. Building upon previous work by \cite{CardShiftShare}, the authors leverage the fact that migrants prefer to settle in areas with an existing stock of migrants from their origin country. However, instead of using the realized ancestry like Card, the authors use only plausibly exogenous variation in pre-existing ancestry as an instrument. In order to do so, they use a simple model of ancestry developed in \cite{MigrantsAncestorsandInvestments}. The model predicts that migration from country $o$ to county $i$ in each period is driven by 3 forces:
    \begin{enumerate}[(i)]
        \item Push Factor: Origin country factors driving emigration from said country to the US,
        \item Pull Factor: Relative attractiveness of a given county to immigrants,
        \item Recursive Factor: Driven by the fact that immigrants from a country prefer moving to places with an existing stock of immigrants from their country.
    \end{enumerate}
        
    The authors then solve this model recursively using the full stock of ancestry data from IPUMS since 1880. They then calculate predicted ancestry from country $o$ to county $d$ at time $t$. The authors argue that this predicted ancestry is plausibly exogenous to both destination-county and origin-country specific factors. They then follow a shift-share approach by interacting predicted ancestry from a country $o$ in county $d$ with the contemporaneous inflow of migrants from said country $o$ into regions of the US other than the one containing $d$ to predict immigration from country $o$ to county $d$ ($\hat{I}_{o,d,t}$). Summing up over all foreign countries, they obtain the predicted immigration $\hat{I}_{d,t}$, the desired instrument. This instrument can then be used to estimate 2SLS coefficients of the effects of immigration on various outcome variables. In relation to my project, the authors estimate and publish immigration shocks for every five-year period from 1975 to 2010 at the county level. I plan to use these immigration shocks as an instrument for Equation \eqref{eq:main}. I present my first-stage results for my instrument in Table \ref{tab:firststage}, across multiple specifications. I henceforth refer to Non-European Migration as Immigration (proxy). The exclusion restriction would be that any confounding factors that drive increases in a given US county’s business dynamics post-1975 do not systematically correlate with pre-1975 immigration from a given origin to other regions within the US, interacted with the simultaneous settlement of European migrants in that US county.

    Table \ref{tab:main} presents my primary results in an IV setting. Column 1 is the log of average annual pay, which can be interpreted as wages, and Column 2 represents the log of average employment, which can also be interpreted as firm size. Column 3 represents the Entry Rate, which measures how many new establishments are created as a percentage of the average number of establishments over a given year. The Exit Rate in column 4 measures the same for establishment death/exit. Since the immigration variable is logged, we can interpret the coefficients as the effects of a 1\% increase in immigration to a US county over 5 years. This will lead to a 0.08 \% increase in annual average wages and a 0.1 \% increase in the average firm size over the 5 years. Simultaneously, we see an increase of 0.183 percentage points in the entry rate and a 0.117 percentage points increase in the exit rate. All the coefficients are statistically significant at the 1 \%. All specifications include state and time fixed effects. The direction of the wage effects is consistent with the finding of \cite{ImmigrationInnovationGrowth}. These effects point to large general equilibrium effects caused by immigration in the county. The findings point to an increase in the business dynamics and wages in a county that receives a plausibly exogenous shock of migrants. I elaborate on these effects in the context of my model in Section \ref{sec:discussion}.
    
    
    For additional robustness, I additionally used a hyperbolic sine on both immigration and the outcome variables. This ensures that the identification approach is robust to any logarithmic issues with zeroes or negative values while ensuring our estimates maintain an elasticity interpretation. The results are presented in Table \ref{tab:sin_iv}, the results remain unchanged. Further, an important point of discussion in the literature concerning business dynamism is population change. Expressed plainly, an increase in population should increase business dynamism, and migration is a form of population growth. To deal with these concerns, in Table \ref{tab:pop}, I control for population by taking the average population of the county over the 5-year period of analysis. All the coefficients remain relatively stable. The results remain unchanged if we use the initial or the final population of the county. This speaks to the robustness of the instruments, along with the fact that the population does not bias my results.

    Finally, in Table \ref{tab:bdsage}, I look for whether there are any heterogeneous effects by age of the firms in a county that receives a positive migration shock. In Panel A, I examine the Exit/Death rate. We see that the increase in exit rate is driven primarily by older firms. It should be noted that since the analysis is performed for 5-year windows, the firms in columns 3 and 4 would have been established before the batch of new migrants predicted by the instrument arrives in the county. It would be plausible to say that this increase in the exit rate isn't caused by new immigrants setting up businesses that fail in short periods. This is not to say that immigrant businesses don't fail at higher rates, but rather that these county-level effects are not primarily driven by new immigrant businesses. There is a larger general equilibrium effect throughout the county. In Panel B, we see that the increase in firm size is spread about equally across the firm age buckets. The estimates for the four age groups are relatively similar and similar to the coefficient for the overall effect in Table \ref{tab:main}. There are no heterogeneous effects by age. I discuss the implications of these findings in the following section.